\section{Peta Konsep: dari Teori Bahasa ke Front-End Kompilator}
\label{sec:ch02-peta}

Konsep-konsep pada bab ini bukan materi hafalan yang berdiri sendiri.
Mereka membentuk rantai yang menghubungkan teori bahasa dengan
implementasi kompilator. Rantai tersebut dapat diringkas sebagai
berikut:
\[
\text{ekspresi reguler} \to \text{NFA} \to \text{DFA} \to
\text{lexer},
\]
\[
\text{CFG} \to \text{derivasi} \to \text{pohon urai} \to
\text{parser} \to \text{AST},
\]
\[
\text{atribut} \to \text{SDD} \to \text{analisis semantik} \to
\text{validasi}.
\]

\begin{figure}[H]
\centering
\begin{tikzpicture}[node distance=6mm and 10mm]
  \node[phase, text width=3.3cm] (teori) {Teori bahasa\\dan otomata};
  \node[phase, text width=3.3cm, right=of teori] (leksikal) {Analisis\\leksikal};
  \node[phase, text width=3.3cm, right=of leksikal] (sintaksis) {Analisis\\sintaksis};
  \node[phase, text width=3.3cm, below=of sintaksis] (semantik) {Analisis\\semantik};
  \node[phase, text width=3.3cm, left=of semantik] (ast) {AST dan\\validasi};
  \draw[arrow] (teori) -- (leksikal);
  \draw[arrow] (leksikal) -- (sintaksis);
  \draw[arrow] (sintaksis) -- (semantik);
  \draw[arrow] (semantik) -- (ast);
\end{tikzpicture}
\caption{Peta konsep dari teori bahasa menuju front-end kompilator}
\label{fig:ch02-peta-konsep}
\end{figure}

\begin{table}[H]
\centering
\caption{Konsep teori dan penerapannya pada front-end}
\label{tab:ch02-peta}
\begin{tabularx}{\textwidth}{@{}l X l@{}}
\toprule
\textbf{Konsep} & \textbf{Penerapan} & \textbf{Bab} \\
\midrule
Ekspresi reguler & Spesifikasi token: identifier, literal, operator &
V \\
NFA dan DFA & Pengenalan token dan konstruksi scanner & VI--VII \\
CFG dan derivasi & Spesifikasi struktur program dan parser & VIII--X \\
FIRST dan FOLLOW & Konstruksi tabel LL(1) & X \\
LR item dan tabel & Konstruksi parser bottom-up & XI--XII \\
Atribut dan SDD & Pemeriksaan tipe dan analisis semantik & XIV \\
Pohon dan AST & Representasi program untuk validasi & VIII, XVI \\
\bottomrule
\end{tabularx}
\end{table}

Peta ini menjelaskan mengapa urutan bab mengikuti urutan fase
kompilasi. Setiap bab menambahkan satu lapisan representasi: dari
karakter menjadi token, dari token menjadi pohon, dari pohon menjadi
struktur bermakna, dan akhirnya menjadi program yang tervalidasi.

Bab III akan memperlihatkan keseluruhan peta ini dalam konteks
arsitektur kompilator. Bab IV kemudian memperdalam hierarki Chomsky,
yang menjadi dasar pemilihan alat pada setiap fase. Dengan memahami
peta konsep ini, mahasiswa dapat melihat bahwa setiap teknik yang
dipelajari memiliki tempat yang jelas dalam sistem yang lebih besar.
